\documentclass[12pt]{article}

\usepackage[utf8]{inputenc} % important for pdflatex + any UTF-8 characters
\usepackage{amsmath, amssymb}
\usepackage{geometry}
\geometry{margin=1in}

\setlength{\parindent}{0pt}

\title{Practice Problem Set 2}
\author{ECON 300 -- Labour Economics}
\date{Winter 2026}

\begin{document}
\maketitle





\subsection*{Problem 1:}

Consider an income maintenance program with the following structure:

\begin{itemize}
\item Guaranteed income (benefit): \( G = 600 \) per week  
\item Benefit reduction rate (tax-back rate): \( t = 0.5 \)  
\item Market wage: \( w = 20 \) per hour  
\end{itemize}

Disposable income is given by:
\[
C = G + (1 - t)wh
\quad \text{for } wh \leq \frac{G}{t}
\]

\begin{enumerate}
\item Calculate the \textbf{break-even level of earnings} where benefits fall to zero.
\item Compute the worker’s \textbf{effective net wage} while receiving benefits.
\item Compare the net wage to the market wage and explain the implication for labour supply incentives.
\item Explain why high benefit reduction rates may discourage work even when market wages are unchanged.
\end{enumerate}

\subsection*{Problem 2: }

A worker has 80 hours per week available for work and leisure.

\begin{enumerate}
\item Draw the budget constraint under:
\begin{itemize}
\item No income support
\item A guaranteed income program with a benefit reduction rate less than 100\%
\end{itemize}
\item Clearly label:
\begin{itemize}
\item The guaranteed income level
\item The break-even point
\end{itemize}
\item Explain how the slope of the budget constraint changes when moving from the no-program case to the program case.
\item Using indifference curves, explain how the program may affect:
\begin{itemize}
\item Labour force participation
\item Hours worked by individuals who remain employed
\end{itemize}
\end{enumerate}

\subsection*{Problem 3: }

Answer each question briefly (3--5 sentences each).

\begin{enumerate}
\item How does a \textbf{negative income tax} differ from a traditional welfare program?
\item Why may an \textbf{earned income tax credit (EITC)} increase labour force participation but reduce hours worked for some individuals?
\item Why do economists distinguish between the \textbf{extensive margin} and the \textbf{intensive margin} of labour supply?
\item Why does targeting income support to low-income workers create trade-offs between equity and efficiency?
\end{enumerate}



\subsection*{Problem 5: }

A competitive firm has the production function:
\[
Q = 10L^{0.5}
\]
where \( L \) is labour.

The firm faces:
\begin{itemize}
\item Output price: \( P = 4 \)
\item Wage rate: \( w = 8 \)
\end{itemize}

\begin{enumerate}
\item Derive the firm’s \textbf{marginal product of labour (MPL)}.
\item Derive the \textbf{value of marginal product of labour (VMPL)}.
\item Determine the firm’s \textbf{optimal employment level}.
\item Explain why the labour demand curve is downward sloping in the wage.
\end{enumerate}

\subsection*{Problem 6:}

\begin{enumerate}
\item Draw the firm’s labour demand curve derived from the VMPL.
\item Show equilibrium employment when:
\begin{itemize}
\item The wage is \( w = 8 \)
\item The wage increases to \( w = 12 \)
\end{itemize}
\item Explain how a change in the wage differs from a change in:
\begin{itemize}
\item Output price
\item Technology
\end{itemize}
\end{enumerate}

\subsection*{Problem 7: }

Using the production function from Problem 5:

\begin{enumerate}
\item Recalculate the firm’s optimal employment level if the output price increases from \( P = 4 \) to \( P = 6 \).
\item Explain the result intuitively using the concept of the value of marginal product.
\item Distinguish clearly between:
\begin{itemize}
\item A movement along the labour demand curve
\item A shift of the labour demand curve
\end{itemize}
\end{enumerate}

\vfill
\textbf{End of Practice Problem Set}

\end{document}
